% ID: SEL-D0029-OB-E02
% ESTADO: validada. S030, cotejada con examen y criterios.
\textbf{a)} En $f$ usamos las reglas del producto y de la cadena.
En $g$ usamos la del cociente y derivamos el logaritmo compuesto.
\[
\begin{aligned}
f'(x)&=e^{5x}(x^2-5)^2[5(x^2-5)+6x],\\
g'(x)&=\frac{6x^2(x^3+1)\ln(x^2+2)-(x^3+1)^2\,2x/(x^2+2)}
{\ln^2(x^2+2)}.
\end{aligned}
\]
El denominador de $g$ no se anula porque $x^2+2\ge2$.

\textbf{b)} La regla del cociente da $h'(x)=-5/(x+5)^2$.
Como $h(0)=2$ y $h'(0)=-1/5$, la tangente es
$\boxed{y=2-x/5}$.
